\cvsection[0.45]{实习经历}

% 公司品牌配置集中放在本模块，便于直接替换颜色、Logo 和名称。
\begingroup
\definecolor{NovaBlue}{HTML}{3C8CFF}
\definecolor{NovaTint}{HTML}{EEF5FF}
\definecolor{HelixRed}{HTML}{CF0A2C}
\definecolor{HelixTint}{HTML}{FFF1F3}

% 可选参数控制招牌上方间距；其余参数依次为背景色、品牌色、
% Logo、公司及部门名称、右侧职位与日期。
\newcommand{\companybanner}[6][0.4em]{%
  \par\nointerlineskip
  \vspace{#1}%
  \begin{tcolorbox}[
    enhanced,colback=#2,frame hidden,boxrule=0pt,
    overlay={%
      \begin{scope}
        \clip[rounded corners=1.5mm]
          (frame.south west) rectangle (frame.north east);
        \fill[#3]
          (frame.south west) rectangle ([xshift=3pt]frame.north west);
      \end{scope}%
      \node[anchor=center,inner sep=0pt]
        at ([xshift=17.25pt]frame.west) {#4};%
    },
    arc=1.5mm,outer arc=1.5mm,boxsep=0pt,
    left=6pt,right=6pt,top=0pt,bottom=0pt,
    height=22pt,valign=center,
    width=\linewidth,before skip=0pt,after skip=0pt]
    \makebox[\linewidth]{%
      \hspace{2.8em}%
      {\fontsize{10.5}{12.6}\selectfont\sffamily\bfseries #5}%
      \hfill
      {\fontsize{10.5}{12.6}\selectfont\sffamily
        \color{OpenCurveMuted}\mbox{#6}}%
    }%
  \end{tcolorbox}%
}

% 主要贡献下的二级要点：首行缩进 0.65em，续行悬挂缩进 1.6em。
\newcommand{\companyitem}[2]{%
  \par\hangindent=1.6em\hangafter=1\noindent
  \hspace{0.65em}\textcolor{#1}{\textbullet}\ #2\par
}

\companybanner{NovaTint}{NovaBlue}{%
  \textcolor{NovaBlue}{\fontsize{18}{18}\selectfont\simpleicon{googlecloud}}%
}{云汇科技\,·\,AI 平台部}{算法工程师实习生 \headerdivider 2025.06--2025.09}
\vspace{0.25em}
\cvtagline{企业知识助手}
\noindent\textbf{项目背景：}面向企业内部知识检索与问答场景，搭建从数据处理、混合检索到答案生成的完整链路。\par
\noindent\textbf{主要贡献：}\par
\companyitem{NovaBlue}{负责核心检索模块与评测流程，支持多种数据源接入。}
\companyitem{NovaBlue}{设计缓存与批处理策略，降低线上响应延迟。}
\noindent\textbf{项目成果：}在内部测试集上将答案准确率从 71\% 提升至 89\%，平均响应时间降低 40\%。\par

\companybanner[0.8em]{HelixTint}{HelixRed}{%
  \textcolor{HelixRed}{\fontsize{18}{18}\selectfont\simpleicon{docker}}%
}{青蓝软件\,·\,研发中心}{软件工程师实习生 \headerdivider 2024.07--2024.10}
\vspace{0.25em}
\cvtagline{研发效能平台}
\noindent\textbf{项目背景：}面向研发流程中的任务编排、结果追踪与重复操作自动化，建设研发效能平台。\par
\noindent\textbf{主要贡献：}\par
\companyitem{HelixRed}{使用 Python 与 PostgreSQL 开发任务编排和结果追踪功能，完善异常处理、日志记录与自动化测试。}
\companyitem{HelixRed}{参与需求评审和代码审查，推动重复操作自动化，将常见任务的人工操作时间减少约 30\%。}

\endgroup
